\chapter{Estrategia 2: swap y collar sobre el WTI}\label{cap:e2}

\begin{objetivos}
\begin{itemize}
\item Describir el diseño de la Estrategia~2: qué se cubre, cuánto y con qué programa rodante.
\item Reproducir los strikes y primas del collar en sus tres capas (Cuadro~7 del informe).
\item Construir la tabla de escenarios sobre el inventario (Cuadro~8) y leer la Figura~5.
\item Medir la cobertura en barriles equivalentes y demostrar que el swap fija el resultado.
\item Evaluar el costo de la backwardation, el skew, el riesgo de base y lo que queda sin cubrir.
\end{itemize}
\end{objetivos}

\begin{prerreq}Capítulos~\ref{cap:futuros} (costo de acarreo, backwardation, base), \ref{cap:swaps} (swap de materias primas), \ref{cap:opciones} (collar) y~\ref{cap:valoracion} (modelo de Black, griegas y skew).\end{prerreq}

\section{El diseño}

\paragraph{La respuesta, primero} La Estrategia~2 (E2) cubre el 80\,\% del inventario de crudo, mitad con un swap y mitad con un collar. Si el WTI cae 30\,\%, la pérdida baja de US\$ 80.6 a 34.0 millones, sin pagar prima. A cambio, si el WTI sube 30\,\%, la ganancia se limita a US\$ 21.8 millones (frente a 80.6 sin cobertura).

\paragraph{Por qué OTC} Se usan contratos con bancos, no futuros de NYMEX, porque los futuros exigen márgenes diarios en efectivo que la liquidez no soporta (cap.~\ref{cap:futuros}).

\subsection{Qué se cubre: un saldo estructural}

El inventario no es un lote fijo: entra y sale crudo todos los días. Con compras de 87 mil barriles diarios, los 2,901 mil barriles en inventario equivalen a $2{,}901/87=33$ días de compras: el inventario \textbf{rota} en aproximadamente un mes.

\begin{definicion}[Saldo estructural]
El \textbf{saldo estructural} es la parte del inventario que, aunque sus barriles cambien, siempre está presente en el tanque: el nivel mínimo con el que opera la refinería. Es un riesgo permanente aunque cada barril individual se venda en semanas.
\end{definicion}

El informe cubre el 80\,\% del inventario: 2.32 millones de barriles, US\$ 214.9 millones al WTI de hoy. El 20\,\% restante queda abierto \emph{a propósito}: si el inventario baja, la empresa no queda sobrecubierta (con más derivados que barriles), lo que sería especulación. El 80\,\% es el tope de la política propuesta (50\,\%--80\,\% del inventario).

\subsection{Cómo se cubre: un programa rodante en tres capas}

\begin{definicion}[Programa rodante]
Un \textbf{programa rodante} cubre un riesgo permanente con contratos de corto plazo escalonados. Cuando vence uno, se contrata otro a plazo más largo, de modo que la cobertura total se mantiene.
\end{definicion}

El volumen cubierto se reparte en \textbf{tres capas} iguales, una por cada mes en que se vende el inventario: noviembre, diciembre y enero. Cada capa es de $2.32/3=0.774$ millones de barriles (774 contratos de 1,000 barriles). Al vencer, cada capa se repone a tres meses. Si no se repusiera, la cobertura caería al 53\,\% del inventario tras el primer vencimiento y al 27\,\% tras el segundo.

Dentro de cada capa, la mitad del volumen va a cada instrumento:
\begin{enumerate}
\item[(a)] Un \textbf{swap de WTI} con liquidaciones mensuales a US\$ 89.38/bl, el promedio de los futuros de noviembre a enero. En total, 1.16 millones de barriles. Se usa para los cargamentos cuyo margen ya está definido: si el precio de venta del producto ya está pactado, conviene fijar también el del crudo.
\item[(b)] Un \textbf{collar de costo cero} por capa: put comprada al 90\,\% del futuro de ese mes (redondeado a 0.10) y call vendida de igual prima. En total, 1.16 millones de barriles.
\end{enumerate}

\begin{intuicion}[por qué mezclar]
El swap elimina la incertidumbre por completo, pero fija hoy un precio inferior al actual por la backwardation y renuncia a toda subida. El collar deja un rango de precios abierto, protege contra caídas grandes y no cuesta prima. Combinándolos, la empresa obtiene un punto intermedio: menos variabilidad que con el collar solo y más posibilidad de beneficiarse de alzas que con el swap solo.
\end{intuicion}

\section{Los parámetros y la valorización del collar}

\begin{center}\small
\captionof{table}{Parámetros de la Estrategia 2 (Cuadro~7 del informe). Las tres cifras van en el orden nov / dic / ene.}\label{tab:c8}
\begin{tabularx}{\linewidth}{@{}>{\raggedright\arraybackslash}Xr@{}}
\toprule
\multicolumn{2}{@{}l}{\emph{Insumos de mercado}}\\
Futuros $F$ (US\$/bl) & 92.60 / 89.01 / 86.54\\
Volatilidad implícita $\sigma$ (OVX) & 56.1\,\%\\
Plazo al vencimiento $T$ (días) & 17 / 50 / 79\\
Tasa libre de riesgo $r$ (SOFR, continua) & 4.12\,\%\\
\midrule
\multicolumn{2}{@{}l}{\emph{Estructura del collar (costo cero, por capa)}}\\
Volumen cubierto: total / en collar (MMbl) & 2.32 / 1.16\\
Put comprada: strike (US\$/bl) & 83.30 / 80.10 / 77.90\\
Put comprada: prima media (US\$/bl) & 3.12\\
Call vendida: strike (US\$/bl) & 103.68 / 100.38 / 98.05\\
\midrule
\multicolumn{2}{@{}l}{\emph{Inventario total si el WTI cae 30\,\% (US\$ MM)}}\\
Sin cobertura & $(80.6)$\\
Solo collar & $(44.4)$\\
\textbf{Propuesta: swap + collar} & $\mathbf{(34.0)}$\\
\bottomrule
\end{tabularx}
\end{center}

Cada capa se valoriza como en el capítulo~\ref{cap:valoracion}: la put al 90\,\% de su futuro ($0.90\times92.60=83.34\to83.30$; $0.90\times89.01=80.11\to80.10$; $0.90\times86.54=77.89\to77.90$) se valoriza con Black, y el strike de la call es el que iguala su prima. Las primas de las puts son 1.12, 3.40 y 4.85: cuanto más plazo, más cara la protección.

\figura{capas}{Piso y techo de cada capa del programa rodante. La banda de cada capa se ensancha con el plazo: con más tiempo hasta el vencimiento, la call de igual prima debe estar más lejos. Los plazos (17, 50 y 79 días) son ilustrativos: al ejecutar se recalculan los strikes con los futuros vigentes.}

\section{El costo de la backwardation}

El swap fija 89.38, por debajo del precio de hoy: frente al spot cede $1.16\times(92.60-89.38)=3.7$ millones (7.5 si todo el 80\,\% fuera swap). Pero, como vimos en el capítulo~\ref{cap:futuros}, ese descuento ya está en el precio futuro: en backwardation (enero 6.5\,\% bajo noviembre) el costo de acarreo implica un rendimiento de conveniencia de 44.7\,\% anual. Si el WTI converge al futuro de enero (86.54), el inventario sin cobertura perdería US\$ 17.6 millones, más que el swap. Frente al futuro, el costo esperado del swap es nulo.

\section{Escenarios}

El informe compara el inventario con y sin cobertura ante \textbf{choques instantáneos} del WTI, aplicados a todas las capas a la vez. Es una prueba de resistencia, no una proyección. Con $Q=2.901$ MMbl en inventario, $Q_C=0.80\,Q=2.32$ cubiertos, $S_0=92.60$ y $S_T$ el WTI del escenario, cada fila aplica una fórmula sencilla. La parte abierta (20\,\%) aporta siempre $(S_T-S_0)\times0.58$:

\begin{center}\small
\begin{tabularx}{\linewidth}{@{}>{\raggedright\arraybackslash}p{0.22\linewidth}X@{}}
\toprule
Posición & Resultado (US\$ millones)\\
\midrule
Sin cobertura & $Q\,(S_T-S_0)$\\
Swap & $Q_C\,(89.38-S_0)+$ abierto\\
Collar & $\frac{Q_C}{3}\sum_{\text{capas}}[\min(\max(S_T,K_p),K_c)-S_0]+$ abierto\\
Solo put & $\frac{Q_C}{3}\sum_{\text{capas}}[\max(S_T,K_p)-S_0-p]+$ abierto\\
Propuesta & $\tfrac12\,\text{Swap}+\tfrac12\,\text{Collar}$\\
\bottomrule
\end{tabularx}
\end{center}

\begin{center}\small
\captionof{table}{Resultado sobre el inventario total, US\$ millones (Cuadro~8 del informe). Coberturas sobre el 80\,\%; Put = solo put con prima pagada; Prop. = 50\,\% swap + 50\,\% collar.}\label{tab:c9}
\begin{tabular}{@{}lrrrrrr@{}}
\toprule
WTI & $S_T$ & Sin cob. & Swap & Collar & Put & Prop.\\
\midrule
$-30\,\%$ & 64.8 & $(80.6)$ & $(23.6)$ & $(44.4)$ & $(51.6)$ & $(34.0)$\\
$-20\,\%$ & 74.1 & $(53.7)$ & $(18.2)$ & $(39.0)$ & $(46.2)$ & $(28.6)$\\
$-10\,\%$ & 83.3 & $(26.9)$ & $(12.8)$ & $(26.9)$ & $(34.1)$ & $(19.9)$\\
$=F$ ene & 86.5 & $(17.6)$ & $(11.0)$ & $(17.6)$ & $(24.8)$ & $(14.3)$\\
$0\,\%$ & 92.6 & 0.0 & $(7.5)$ & 0.0 & $(7.2)$ & $(3.7)$\\
$+10\,\%$ & 101.9 & 26.9 & $(2.1)$ & 22.8 & 19.6 & 10.3\\
$+20\,\%$ & 111.1 & 53.7 & 3.3 & 29.6 & 46.5 & 16.4\\
$+30\,\%$ & 120.4 & 80.6 & 8.7 & 34.9 & 73.3 & 21.8\\
\bottomrule
\end{tabular}
\end{center}

\begin{ejemplo}{Tres celdas del cuadro}
\emph{Collar, WTI $-30\,\%$} ($S_T=64.82$): todas las capas están bajo su piso. Parte cubierta: $0.774\times[(83.30-92.60)+(80.10-92.60)+(77.90-92.60)]=0.774\times(-36.5)=-28.2$. Parte abierta: $0.580\times(64.82-92.60)=-16.1$. Total: $-44.4$. Observe que la parte abierta no tiene piso: por eso el collar no protege \enquote{al 100\,\%}.

\emph{Propuesta, WTI $-10\,\%$} ($S_T=83.34$): swap $2.32\times(89.38-92.60)+0.580\times(83.34-92.60)=-7.5-5.4=-12.8$. Collar: el precio está entre el piso y el techo de todas las capas (apenas sobre el piso de noviembre, 83.30), así que no hay protección: $2.32\times(-9.26)-5.4=-26.9$. Promedio: $(-12.8-26.9)/2=-19.9$.

\emph{Solo put, WTI $-30\,\%$}: $0.774\times[(83.30-92.60-1.12)+(80.10-92.60-3.40)+(77.90-92.60-4.85)]-16.1=0.774\times(-45.87)-16.1=-51.6$.
\end{ejemplo}

\figura{e2_escenarios}{Perfil de resultados de las alternativas para el inventario (amplía la Figura~5 del informe con la put sola). Con el 20\,\% abierto ninguna línea es plana: hasta el swap pierde si el WTI cae. La propuesta (azul grueso) queda entre el swap y el collar.}

\section{Evaluación}

\begin{center}\small
\captionof{table}{Resumen de las alternativas (US\$ millones), a partir del Cuadro~8 del informe.}
\begin{tabularx}{\linewidth}{@{}>{\raggedright\arraybackslash}Xrrr@{}}
\toprule
Alternativa (sobre el 80\,\%) & Prima & WTI $-30\,\%$ & WTI $+30\,\%$\\
\midrule
Sin cobertura & 0 & $(80.6)$ & 80.6\\
Solo swap & 0 & $(23.6)$ & 8.7\\
Solo collar & 0 & $(44.4)$ & 34.9\\
Solo put & 7.2 & $(51.6)$ & 73.3\\
50\,\% swap + 50\,\% collar & 0 & $(34.0)$ & 21.8\\
\bottomrule
\end{tabularx}
\end{center}

\paragraph{Por qué 50/50} El reparto equilibra los extremos. Solo con swap, la pérdida en el peor escenario sería de 23.6 millones, pero con un alza de 30\,\% la ganancia se reduciría a 8.7. Solo con collar serían 44.4 y 34.9. La mezcla da 34.0 y 21.8.

\paragraph{Por qué no la put sola} Tiene el mejor perfil (piso y casi toda la subida), pero cuesta 7.2 millones de prima por cada rotación del programa. El informe la deja como plan~B para el tramo más expuesto, porque no exige colateral.

\paragraph{Riesgo residual} El VaR al 95\,\% al vencimiento de las opciones baja de US\$ 93.3 a 36.5 millones: una reducción de 61\,\%, menor que en la E1 porque queda abierto el 20\,\% del inventario y la zona entre los strikes del collar.

\section{La cobertura medida en barriles}

Así como la E1 mide posiciones de cambio (PCC, PND, PCG), la E2 mide su posición en barriles equivalentes:

\begin{center}\small
\captionof{table}{Posiciones en barriles (Cuadro~9 del informe). Positivo: largo en crudo.}
\begin{tabularx}{\linewidth}{@{}>{\raggedright\arraybackslash}Xrr@{}}
\toprule
& MMbl & WTI $-10\,\%$ (US\$ MM)\\
\midrule
Físico: inventario de crudo & 2.90 & $(26.86)$\\
Swap (vende a precio fijo) & $(1.16)$ & 10.75\\
Collar (equivalente delta) & $(0.62)$ & 5.75\\
\midrule
Global (físico + derivados) & 1.12 & $(10.37)$\\
\bottomrule
\end{tabularx}
\end{center}

El swap es lineal: equivale a vender 1.16 millones de barriles. El collar equivale a vender $0.535\times1.16=0.62$ millones: su delta media por barril es $-0.535$ (cap.~\ref{cap:valoracion}). La posición larga baja de 2.90 a 1.12 millones de barriles, un 61\,\% menos. Lo que queda es el 20\,\% abierto y la zona entre los strikes, donde el collar no protege.

\paragraph{Demostración de la cobertura} En la parte que cubre el swap, lo que pierde el crudo lo gana el swap (Cuadro~10 del informe):
\begin{center}\small
\begin{tabular}{@{}lrrr@{}}
\toprule
& WTI $-30\,\%$ & $=F$ ene & WTI $+30\,\%$\\
\midrule
Inventario (1.16 MMbl) & $(32.2)$ & $(7.0)$ & 32.2\\
Swap: $(89.38-S)\times1.16$ & 28.5 & 3.3 & $(36.0)$\\
\midrule
Total & $(3.7)$ & $(3.7)$ & $(3.7)$\\
\bottomrule
\end{tabular}
\end{center}
El total es constante: $(89.38-92.60)\times1.16=-3.7$, el descuento de la backwardation.

\section{Sensibilidades y riesgo de base}

\paragraph{Volatilidad} El OVX (56.1\,\%) coincide con la volatilidad realizada del futuro en el último año (55.8\,\%). El collar casi no depende del nivel de volatilidad: pierde unos 9 mil dólares por punto (cap.~\ref{cap:valoracion}).

\paragraph{Skew} Sí depende del skew: si las calls cotizaran 5 puntos de volatilidad bajo las puts, el techo de enero bajaría de 98.05 a 95.37; con 10 puntos, a 92.90. Por eso los strikes se fijan con cotizaciones bancarias, no con el modelo.

\paragraph{Base} El riesgo residual principal es la base: se compran crudos referidos al Brent, hoy US\$ 12.7/bl sobre el WTI frente a un promedio de 4.1 en dos años. El ratio de mínima varianza del Brent frente al futuro de WTI es 1.14, con una efectividad de solo 79\,\% (cap.~\ref{cap:futuros}). Como los escenarios suponen que el inventario sigue al WTI, \emph{sobrestiman} la reducción del riesgo. Se mantiene el WTI por su liquidez y porque los precios de venta siguen la paridad de importación sobre el WTI.

\begin{cuidado}[descalce de fechas]
Las opciones vencen antes que su futuro (7 días hábiles antes del 26 del mes previo), el swap se liquida mes a mes y el inventario se vende a lo largo de cada mes. Las fechas no coinciden exactamente: parte del riesgo de base es temporal, y el programa rodante debe ajustarse al calendario real de ventas.
\end{cuidado}

\section{Lo que la E2 no cubre: el margen}

La E2 fija el precio del crudo en inventario, no el margen de refinación. Quedan fuera el inventario de productos (US\$ 31.3 millones por cada 10\,\% de caída) y la variación del \emph{crack} entre compra y venta. Para ellos existen swaps líquidos de diésel y gasolina (NYMEX HO y RBOB); el informe los deja como fase~2 (cap.~\ref{cap:diagnostico}).

\section{Valor razonable: un adelanto}

El swap y el collar nacen con valor cero. Si la curva de futuros cayera 10\,\% de inmediato, el collar sería un activo de US\$ 5.7 millones y el swap uno de 10.3 (8.86 por barril por 1.16 millones de barriles); con un alza de 10\,\%, serían pasivos de 5.8 y 10.3 millones. El capítulo~\ref{cap:niif} explica cómo se registran.

\begin{resumen}
\begin{itemize}
\item El inventario rota en unos 33 días, pero tiene un saldo estructural. La E2 cubre el 80\,\% (2.32 MMbl, US\$ 214.9 millones) con un programa rodante de tres capas (nov, dic, ene) de 0.774 MMbl cada una; el 20\,\% queda abierto para no sobrecubrir.
\item Mitad swap a 89.38 con liquidaciones mensuales; mitad collar de costo cero por capa: pisos 83.30 / 80.10 / 77.90 y techos 103.68 / 100.38 / 98.05.
\item Ante una caída de 30\,\% del WTI, la pérdida baja de 80.6 a 34.0 millones sin prima; con un alza de 30\,\%, la ganancia se limita a 21.8. El VaR al 95\,\% baja de 93.3 a 36.5.
\item La posición larga en crudo baja de 2.90 a 1.12 MMbl equivalentes (61\,\%).
\item El swap cede 3.7 millones frente al spot, pero nada frente al futuro. El collar depende poco de la volatilidad y mucho del skew.
\item Los riesgos residuales son la base Brent--WTI, el 20\,\% abierto, la zona entre strikes y el margen de refinación (fase~2).
\end{itemize}
\end{resumen}

\section*{Ejercicios}
\addcontentsline{toc}{section}{Ejercicios}

\ejercicio[cálculo] Con el modelo del Cuadro~8, calcule el resultado del swap, del collar y de la propuesta si el WTI termina en 95.00 y en 70.00.

\ejercicio[cálculo] Suponga que la empresa repartiera lo cubierto en 70\,\% swap y 30\,\% collar. Calcule el resultado con el WTI $-30\,\%$ y $+30\,\%$.

\ejercicio[cálculo] Verifique la celda \enquote{Solo put, $=F$ ene} del Cuadro~8 del informe.

\ejercicio[conceptual] ¿Por qué tiene sentido usar el swap para los cargamentos \enquote{cuyo margen ya está definido} y el collar para el resto?

\ejercicio[conceptual] Si la backwardation desapareciera y la curva quedara plana en 92.60, ¿cómo cambiaría el resultado del swap? ¿Seguiría siendo útil el collar?

\ejercicio[conceptual] ¿Por qué el informe descarta comprar solo puts como estrategia principal, pero las conserva como plan~B?

\ejercicio[cálculo] Vence la capa de noviembre y, por falta de líneas con los bancos, no se repone. ¿Qué porcentaje del inventario queda cubierto? ¿Y si tampoco se repone la de diciembre?
